\section{\boxbox{}: Method Details}
\label{app:boxbox}

Figure~\ref{fig:boxbox-pseudocode} summarizes the task and review loops.
The following explains how harness revisions affect subsequent interactions
and how review, hot reloading, and recovery preserve task continuity.

\begin{figure}[H]
  \centering
\begingroup
\small
\definecolor{heading}{HTML}{385E8D}
\setlength{\parindent}{0pt}
\newenvironment{pseudocode}{%
  \begin{list}{\color{gray}\scriptsize\arabic{enumi}}{%
    \usecounter{enumi}
    \setlength{\leftmargin}{1.7em}
    \setlength{\labelwidth}{1.2em}
    \setlength{\labelsep}{.5em}
    \setlength{\itemsep}{1.5pt}
    \setlength{\parsep}{0pt}
    \setlength{\topsep}{4pt}}%
}{\end{list}}
\newcommand{\Line}[2]{\item\hspace*{#1em}#2}
\newcommand{\kw}[1]{\textbf{#1}}
\begin{minipage}{\linewidth}
{\normalsize\bfseries\color{heading} \boxbox{}}
\par\vspace{5pt}\hrule\vspace{6pt}
\textbf{State:} harness $H$, runtime $R$, task environment $e$, task trajectory $\tau$.
\par\textbf{Review state:} environment $e_r$ and separate trajectory $\tau_r$.
\par\textbf{Budget:} shared task budget $B$.

\begin{pseudocode}
  \Line{0}{\kw{while} budget $B$ remains \kw{do}}
  \Line{1}{$a \sim p_{\mathcal{M},H}(\cdot\mid\tau)$}
  \Line{1}{\kw{if} $a=\texttt{boxbox()}$ \kw{then}}
  \Line{2}{Yield control to the kernel; suspend task execution and preserve context}
  \Line{2}{Prepare $(e_r,\tau_r)$ for independent review under $H$, with access to $\tau$}
  \Line{2}{\kw{while true do}}
  \Line{3}{$a_r \sim p_{\mathcal{M},H}(\cdot\mid\tau_r)$}
  \Line{3}{\kw{if} $a_r=\texttt{commit}(H')$ \kw{then break}}
  \Line{3}{$(e_r,o_r) \gets \operatorname{Execute}_H(e_r,a_r)$}
  \Line{3}{$\tau_r \gets \tau_r\oplus(a_r,o_r)$}
  \Line{2}{\kw{end while}}
  \Line{2}{\kw{if} $H'\ne H$ \kw{then}}
  \Line{3}{The kernel validates $H'$ and attempts activation of its runtime $R'$}
  \Line{3}{\kw{if} successful \kw{then} select $(H,R)\gets(H',R')$}
  \Line{3}{\kw{else} retain or restore the previous harness and runtime}
  \Line{2}{\kw{end if}}
  \Line{2}{Restore continuation context and return control to the selected runtime}
  \Line{1}{\kw{else}}
  \Line{2}{$(e,o) \gets \operatorname{Execute}_H(e,a)$}
  \Line{2}{$\tau \gets \tau\oplus(a,o)$}
  \Line{1}{\kw{end if}}
  \Line{0}{\kw{end while}}
\end{pseudocode}
\vspace{4pt}\hrule
\end{minipage}
\endgroup
  \caption{\textbf{\boxbox{} pseudocode.} Task execution and independent harness review form nested interaction loops, with kernel-managed hot reloading connecting them.}
  \label{fig:boxbox-pseudocode}
\end{figure}

\subsection{Effects on Subsequent Interactions}
\label{app:boxbox-update-effects}

Equation~\ref{eq:boxbox-interaction} describes two ways a harness revision
can affect subsequent execution while the language model $\mathcal{M}$ remains
fixed. Changes to how the harness constructs model inputs can alter
$p_{\mathcal{M},H_k}$, influencing the agent's next action. Changes to
$\operatorname{Execute}_{H_k}$ can alter how that action is carried out and how
its result is returned as an observation. A revision can affect either or
both, so its influence extends to the trajectory generated after activation.

The current harness also governs the review loop. Once a revision is
activated, subsequent reviews use the updated implementation and can examine
the experience accumulated under it. Harness changes can therefore affect
both task execution and the process used to produce later revisions.
Appendix~\ref{app:case-study} illustrates concrete changes to guidance,
measurement, command polling, and context handling.

\subsection{Review}
\label{app:boxbox-agent-initiated}

The task-solving agent requests review through \texttt{boxbox()} without
having to propose a patch. The initial harness requires intervening task
activity and a 30-minute cooldown after each completed review; once eligible,
the agent decides whether and when to request another review.
The kernel preserves the task conversation and launches a separate agent
review session under the current harness. The reviewer can inspect task
trajectories and Git history, but edits only the harness, treating project
code as read-only. It may retain the current harness, add guidance, or modify
executable behavior, committing any changes before ending review.
The default review skill is provided in Appendix~\ref{app:boxbox-review-skill}.

\subsection{Execution Context}
\label{app:boxbox-execution-context}

Execution context includes both runtime-owned state and the existing task
environment. The former includes the task conversation, review history, and
records of which trajectory entries have already been reviewed. These are
exported as a state snapshot. The terminal session and project workspace
remain in place: shell state, ongoing processes, and project files are
preserved through the existing environment connection.

When the agent calls \texttt{boxbox()}, the runtime yields control to the
kernel, suspending the task interaction loop while review proceeds separately.
This pauses the agent's task decisions, not the entire container or its
processes. Before importing state into a candidate runtime, the kernel makes
a separate copy, preserving the original snapshot for recovery if the
candidate fails. This separates changing the agent implementation from
restarting its task environment.

\subsection{Hot Reloading}
\label{app:boxbox-transition}

Harness revisions are tracked as Git commits, while the kernel remains
separate from the editable harness. Before switching, the kernel checks the
committed candidate, loads its runtime, and verifies that it supports task
execution and context transfer. The current runtime keeps its terminal
connection until the candidate successfully imports the copied state.
The kernel then transfers the same terminal session to the candidate and
returns control to it, resuming task interaction from the preserved
conversation in the same workspace. If activation fails, the kernel retains
or restores the previous runtime and context.
For agent-requested reviews, the reviewer can leave concise guidance on using
new capabilities and continuing the task. This guidance is preserved across
the update, and the resumed agent is prompted to read it before continuing.
Task execution and review share the same budget and cost accounting.
Table~\ref{tab:boxbox-continuation} summarizes how execution continues under
each outcome.

\begin{table}[htbp]
  \centering
  \small
  \caption{Task continuation across review and runtime-update outcomes.}
  \label{tab:boxbox-continuation}
  \renewcommand{\arraystretch}{1.15}
  \begin{tabularx}{\linewidth}{@{}>{\raggedright\arraybackslash}p{.26\linewidth}X@{}}
    \toprule
    Outcome & Continuation \\
    \midrule
    No harness change &
    Resume the current runtime with the preserved task context. \\
    Successful update &
    Continue under the revised harness with transferred context, the existing
    terminal session, and the same project workspace. \\
    Update fails &
    Retain or restore the previous runtime and its preserved context. \\
    Runtime fault &
    Enter recovery and review using the retained execution records; restart
    the terminal while preserving project files. \\
    \bottomrule
  \end{tabularx}
\end{table}

\subsection{Recovery}
\label{app:boxbox-validation}

The kernel retains a recovery point across harness updates. A newly activated
runtime remains provisional until it reaches the next agent-requested review
without a runtime fault. If it fails before that boundary, the kernel restores
the previous confirmed harness and state and initiates a review using the
fault evidence. Faults in a confirmed runtime also trigger review.
Runtime-fault recovery starts a fresh terminal, so shell-local state is lost,
but project files, harness history, and execution records remain available.
The retained records let subsequent reviews examine failed revisions and help
the task-solving agent determine how to continue interrupted work.
